\section{An empirical norm comparison for functions with bounded slope variation}
\label{sec:sharp-design}
Assume $0<p_-\le p(x)\le p_+<\infty$ on $[-1,1]$.
For each fixed $0<\eta\le\eta_0$, and for continuous $f$ with finite distributional second-derivative
variation $V=\|f''\|_{\rm TV}$, there is a common event on which
\[
 |P_nf^2-P_Xf^2|
 \le C\eta P_Xf^2+
 C\eta^{-5}\left(\frac{\log(Cn/\delta)}n\right)^3V^2,
 \qquad 0<\eta\le\eta_0, \tag{S}
\]
provided $n\eta^2\ge C\log(Cn/\delta)$.
The event is uniform over all such $f$, with no bound on its affine
part. Constants depend only on $p_-,p_+$.

\begin{proof}
Use $m$ equal bins, $\Delta=2/m$, the interpolant $s=I_m f$, and
$r=f-s$. Let $t=C\log(Cn/\delta)$ and choose
$m=\lfloor c n\eta^2/t\rfloor\ge1$. The usual matrix-Chernoff event
for the hat-function space gives
\[
 |(P_n-P_X)s^2|\le\eta P_Xs^2. \tag{S1}
\]
For completeness, the population Gram matrix of the $m+1$ linear hat
functions has all eigenvalues between $c/m$ and $C/m$: integrate the
square of a linear interpolant on each bin, then use the two density
bounds. The vector of hat values has at most two nonzero entries, each
at most one. After population whitening its squared norm is at most
$Cm$. Matrix Chernoff therefore gives (S1) when
$n\ge C m\eta^{-2}\log(Cm/\delta)$, as ensured by the chosen $m$.
The population Gram bound also gives
\[
 \max_j|s(t_j)|\le C\sqrt m\,\|s\|_{L^2(P_X)}. \tag{S2}
\]

For a bin $I=[a,b]$, define the zero-boundary Green function
\[
 G_u(x)=(x-u)_+
       -\frac{b-u}{b-a}(x-a),\quad x\in I,\ u\in I,
\]
and set it to zero outside $I$. Then
$r(x)=\int_I G_u(x)\,\mu(du)$ on $I$, where $\mu=f''$.
The sign of this formula is immaterial. We have
$|G_u|\le\Delta$, and $u\mapsto G_u$ is uniformly $2$-Lipschitz.
Atoms at grid endpoints have zero Green function and do not affect
the residual. Let $\psi_{I,0},\psi_{I,1}$ be the two affine endpoint
hat functions restricted to $I$.

Bernstein's inequality and a net in $u$ give, simultaneously over all
bins, their two endpoint functions and all $u\in I$,
\[
 |(P_n-P_X)(\psi_{I,j}G_u)|
 \le C\{\Delta^{3/2}\sqrt{t/n}+\Delta t/n\}. \tag{S3}
\]
Indeed the summands have magnitude at most $\Delta$ and variance
at most $C\Delta^3$. A net of spacing $\Delta n^{-10}$ has polynomial
size in $n$ in each bin, and the uniform Lipschitz bound controls
interpolation from the net. The negligible net error can be
absorbed in the displayed right side.

Similarly, a two-dimensional net in $u,v$ gives
\[
 |(P_n-P_X)(G_uG_v)|
 \le C\{\Delta^{5/2}\sqrt{t/n}+\Delta^2t/n\}
 \quad(u,v\in I), \tag{S4}
\]
because the magnitude is at most $\Delta^2$, the variance is at
most $C\Delta^5$, and the parameter Lipschitz bound is at most
$C\Delta$. Kernels in different bins have disjoint supports.
The union of (S1), (S3), and (S4) has probability at least $1-\delta$
after choosing constants in $t$ and in the sample-size condition.

Let $V_I=|\mu|(I^\circ)$, so $\sum_I V_I\le V$.
Equations (S2)--(S3) imply
\[
 |(P_n-P_X)(sr)|
 \le C V\|s\|_{L^2(P_X)}
       \left\{m^{-1}\sqrt{t/n}+m^{-1/2}t/n\right\}. \tag{S5}
\]
The two endpoint coefficients on a bin are bounded by the left
side of (S2); summing the signed-measure integrals therefore gives
(S5). Equation (S4) and $\sum_I V_I^2\le V^2$ give
\[
 |(P_n-P_X)r^2|
 \le C V^2\{m^{-5/2}\sqrt{t/n}+m^{-2}t/n\}. \tag{S6}
\]
The deterministic interpolation bound is
\[
 P_Xr^2\le C V^2m^{-3}. \tag{S7}
\]

Put $a=m^{-1}\sqrt{t/n}+m^{-1/2}t/n$.
By Young's inequality, (S5) is bounded by
$\eta P_Xs^2+C\eta^{-1}V^2 a^2$.
Since $mt/n\le c\eta^2$,
\[
 \eta^{-1}a^2\le C\eta m^{-3},\qquad
 m^{-5/2}\sqrt{t/n}+m^{-2}t/n\le C\eta m^{-3}.
\]
Thus (S1), (S5), and (S6) yield
\[
 |(P_n-P_X)f^2|
 \le C\eta P_Xs^2+C\eta V^2m^{-3}.
\]
Finally $P_Xs^2\le2P_Xf^2+2P_Xr^2$, so (S7) proves the same bound
with $f$ in place of $s$. Substituting $m\asymp n\eta^2/t$ gives
the remainder $C\eta^{-5}(t/n)^3V^2$ in (S).
\end{proof}
